\subsection{Greedy decoding: sequence-level and training errors}

Training does not control errors on greedy paths. Off-policy distillation controls per-token errors on
teacher-sampled prefixes, and on-policy distillation controls them on student-sampled prefixes. Both are
tied to sequence-level divergences by \cref{thm:chain,prop:hybrid}. A teacher sample passes through
$s^\star_t$ only with probability $\pi^\star_t$, and this turns out to be the exact exchange rate.

\begin{theorem}[Path-weighted thresholds]\label{thm:greedy-seq}
\Pv\Ck For every $f$-divergence and every input $x$,
\[
\inf_{q_\theta}\Big\{D_f(\Pseq\|\Qseq):\ g_{q_\theta}(x)\ne g_p(x)\Big\}=\omega_f(x):=\min_{1\le t\le T}k_f\big(\pi^\star_tp^\star_{t,1},\,\pi^\star_tp^\star_{t,2}\big).
\]
In particular, writing $\omega_{\mathrm{rKL}}$ for the threshold of $\KL(\Qseq\|\Pseq)$,
\begin{gather*}
\omega_{\TV}=\tfrac12\min_t\pi^\star_t\gamma_t,\qquad
\omega_{\KL}=\min_t\pi^\star_t\,\kappa_{\KL}\big(p_t(\cdot\mid s^\star_t)\big),\\
\omega_{\mathrm{rKL}}=\min_t\,-\log\Big(1-\pi^\star_t\big(\sqrt{p^\star_{t,1}}-\sqrt{p^\star_{t,2}}\big)^2\Big).
\end{gather*}
Consequently $\Ldec(\theta;\mu)\le\Phi_{\omega_f}\big(\cL_{D_f}(\theta;\mu)\big)\le\bar\Phi_{\omega_f}\big(\cL_{D_f}(\theta;\mu)\big)$. The
supremum of $\Ldec$ over students with $\cL_{D_f}\le\eps$ obeys the two-sided bound of \cref{thm:greedy}(iii),
with $\omega_f$ in place of $\bar\kappa$. Under $\mu=\cD$, \cref{thm:chain} gives $\cL_{\KL}=T\eoff^{\KL}$ and
$\cL_{\mathrm{rKL}}=T\eon^{\mathrm{rKL}}$, and \cref{prop:hybrid} gives $\cL_{\TV}\le T\min\{\eoff^{\TV},\eon^{\TV}\}$. For TV and
KL, the extremal students change one next-token distribution per input, and for them all these relations
are equalities. So the bounds in terms of per-token training errors are sharp too.
\end{theorem}
\begin{proof}
Fix $x$ and drop it from the notation.
\emph{Lower bound.} Let $t$ be the first flip (\cref{lem:firstflip}). Up to step $t$ the student follows the
teacher's greedy path, so $Q(s^\star_t)=\prod_{j<t}q_j(a^\star_j\mid s^\star_j)>0$. The flip then gives
$Q(s^\star_t,a^\star_t)\le Q(s^\star_t,b)$ for some $b\ne a^\star_t$. Apply \cref{lem:merge} to the law of $y_{\le t}$ with
$S_+=\{(s^\star_t,a^\star_t)\}$ and $S_-=\{(s^\star_t,b):b\ne a^\star_t\}$, whose teacher masses are $\pi^\star_tp^\star_{t,1}$ and at most
$\pi^\star_tp^\star_{t,2}$. Finally, $D_f(P\|Q)\ge D_f(P_{\le t}\|Q_{\le t})$ by data processing.
\emph{Attainment.} We may assume $\omega_f(x)<\infty$. Choose a minimising $t$ with $p^\star_{t,2}>0$ if $f(0)=\infty$.
This is possible. If $p^\star_{t,2}=0$, let $j<t$ be the last step with $p^\star_{j,2}>0$. The steps strictly
between $j$ and $t$ are deterministic, so $\pi^\star_jp^\star_{j,1}=\pi^\star_t$. Apply \cref{lem:merge} to a law with atoms
of masses $\pi^\star_t$, $P_2:=\pi^\star_jp^\star_{j,2}$ and $0$, first with $S_-$ the null atom and then with $S_-$ both other
atoms. The second constraint set is larger, so $k_f(\pi^\star_t,P_2)\le k_f(\pi^\star_t,0)$, and $j$ is also a
minimiser. Such a $j$ exists, since otherwise $\pi^\star_t=1$ and $k_f(1,0)=\infty$. At this $t$, take the law of \cref{lem:merge} for $y_{\le t}$ and extend it
to sequences with the teacher's conditional law given $y_{\le t}$. Its likelihood ratio is constant on the three
cells, so $D_f(P\|Q)=\omega_f(x)$. With the strict version of \cref{lem:merge}, its greedy path either leaves
the teacher's before step $t$ or flips at $s^\star_t$. Either way the outputs differ, at cost at most
$\omega_f(x)+\delta$. States of different inputs are disjoint, so the per-input constructions combine. The
probability statements then follow as in the proof of \cref{thm:greedy}(ii)--(iii).
\emph{Special cases.} For TV and KL, the minimiser in $k_f$ is $c=\frac{P_1+P_2}2$ and leaves the rest unchanged.
So $k_f(\pi P_1,\pi P_2)=\pi\,k_f(P_1,P_2)$, and the extremal law changes only $q_t(\cdot\mid s^\star_t)$. Both prefix laws
then put mass $\pi^\star_t$ on $s^\star_t$ and agree elsewhere, so the hybrid sums of \cref{prop:hybrid} equal
$\TV(P,Q)$. For reverse KL, use the closed form of \cref{prop:flip}(ii) with $(p_1,p_2)$ replaced by
$\pi^\star_t(p^\star_{t,1},p^\star_{t,2})$.
\end{proof}

\begin{corollary}[Behavioural and decoding losslessness are independent]\label{cor:separation}
\Pv\Ck
\begin{enumerate}[label=(\roman*),nosep]
\item (Confident teachers.) If $\pi^\star_T(x)\ge\pi_0$ for $\mu$-almost every $x$, then $\omega_{\TV}\ge\pi_0\bar\kappa_{\TV}$ and
$\omega_{\KL}\ge\pi_0\bar\kappa_{\KL}$. Sequence-level losslessness then controls decoding as in \cref{thm:greedy}, at
the price of dividing the budget by $\pi_0$.
\item (Uncertain teachers.) Let the teacher emit i.i.d.\ tokens from $\{0,1\}$ with $p(0)=\frac{1+\gamma}2$. Change only the
student's next-token distribution at the prefix $0^{T-1}$, to $(\frac12-\delta,\frac12+\delta)$. Then the greedy output of
every input changes, while $\TV(\Pseq,\Qseq)=(\frac\gamma2+\delta)(\frac{1+\gamma}2)^{T-1}$ and
$\KL(\Pseq\|\Qseq)\to(\frac{1+\gamma}2)^{T-1}\kappa_{\KL}$ as $\delta\to0$. For $\gamma=0.2$ and $T=100$ both are below
$2\cdot10^{-23}$, yet $\Ldec=1$.
\item (Converse.) The student that emits $\gr(p_t(\cdot\mid s))$ with probability one at every state has $\Ldec=0$. Yet
$\TV(\Pseq,\Qseq)=1-\Pseq(g_p(x))$, and $\KL(\Pseq\|\Qseq)=\infty$ unless $\Pseq$ is a point mass.
\end{enumerate}
The path weight $\pi^\star_t$ is thus the exact exchange rate between the two notions. It is close to $1$ only
when the teacher is confident along its own greedy path.
\end{corollary}
\begin{proof}
(i) $\pi^\star_t\ge\pi^\star_T$ for $t\le T$, and $k_f(\pi P_1,\pi P_2)=\pi\,k_f(P_1,P_2)$ for TV and KL.
(ii) The two laws differ only in the next-token law at $s=0^{T-1}$, which the teacher reaches with probability
$(\frac{1+\gamma}2)^{T-1}$. When two autoregressive laws differ at a single state $s$,
$\TV(\Pseq,\Qseq)=\Pseq(s)\,\TV(p(\cdot\mid s),q(\cdot\mid s))$. The KL value follows from \cref{thm:chain}.
(iii) $\Qseq$ is the point mass at $g_p(x)$.
\end{proof}

\begin{corollary}[Rates under a margin condition]\label{cor:rates}
\Pv\Ck Let $\bar\gamma(x):=\min_t\pi^\star_t(x)\gamma_t(x)$ be the path-weighted greedy margin, and suppose that
$\mu[\bar\gamma\le u]\le Cu^\alpha$ for all $u>0$. This is a margin condition of the type introduced by
\citet{mammen1999smooth} and \citet{tsybakov2004optimal}. Then for every student
\[
\Ldec\le\Big(\frac{2(1+\alpha)}{\alpha}\,\cL_{\TV}\Big)^{\frac\alpha{1+\alpha}}C^{\frac1{1+\alpha}}
\qquad\text{and}\qquad
\Ldec\le\Big(\frac{2(2+\alpha)}{\alpha}\,\cL_{\KL}\Big)^{\frac\alpha{2+\alpha}}C^{\frac2{2+\alpha}}.
\]
Both exponents are attained. Take $T=1$, $V=2$, and $\mu$ uniform on $n$ inputs with margins $(i/n)^{1/\alpha}$,
$i\le n$, so that $C=1$. As $n\to\infty$, the worst-case $\Ldec$ at $\cL_{\TV}=\eps\le\frac{\alpha}{2(1+\alpha)}$ converges to
the first bound; for larger $\eps$ the bound exceeds $1$.
At $\cL_{\KL}=\eps$, its ratio to the second bound tends to $1$ as $n\to\infty$ and then $\eps\to0$. Through
Pinsker's inequality, the first bound would give only the exponent $\frac\alpha{2(1+\alpha)}$ in $\cL_{\KL}$. If instead
the per-token TV on greedy paths is at most $\eta$ everywhere and $\mu[\min_t\gamma_t\le u]\le C_\star u^\alpha$, then
$\Ldec\le C_\star(2\eta)^\alpha$ by \cref{thm:greedy}(iv).
\end{corollary}
\begin{proof}
By \cref{thm:greedy-seq,lem:knapsack}, $\Ldec\le\bar\Phi_\omega(\eps)$ with $\eps=\cL_{D}$. Moreover
$\E(1-\omega/\tau)_+=\frac1\tau\int_0^\tau\mu[\omega<u]\dd u$. If $\mu[\omega<u]\le C'u^a$ for all $u$, then
\[
\bar\Phi_\omega(\eps)\le\inf_{\tau>0}\Big\{\frac\eps\tau+\frac{C'\tau^a}{1+a}\Big\}=\Big(\frac{(1+a)\eps}{a}\Big)^{\frac a{1+a}}C'^{\frac1{1+a}},
\]
with the infimum at $\tau^{1+a}=\frac{(1+a)\eps}{aC'}$. For TV, $\omega_{\TV}=\bar\gamma/2$, so $(a,C')=(\alpha,2^\alpha C)$. For KL,
$\pi^\star_t\kappa_{\KL}\ge\pi^\star_t\gamma_t^2/(2m_t)\ge(\pi^\star_t\gamma_t)^2/2$ by \cref{prop:flip}(ii) and $\pi^\star_tm_t\le1$. So
$\omega_{\KL}\ge\bar\gamma^2/2$ and $(a,C')=(\alpha/2,2^{\alpha/2}C)$. In the example, the function
$\tau\mapsto\frac\eps\tau+\frac1\tau\int_0^\tau\mu[\omega<u]\dd u$ has derivative
$\tau^{-2}\big[\int_0^\tau(\mu[\omega<\tau]-\mu[\omega<u])\dd u-\eps\big]$, whose bracket is non-decreasing. So the function is
quasi-convex, and its stationary point is its minimum. In the TV case the limit law satisfies
$\mu[\omega<u]=(2u)^\alpha$ for $u\le\frac12$. The stationary point satisfies $\tau\le\frac12$ exactly when
$\eps\le\frac\alpha{2(1+\alpha)}$, and then the display is an equality for the limit law. The discrete law converges to it,
and \cref{lem:knapsack} loses at most $1/n$. For
$T=1$ and $V=2$, $\kappa_{\KL}=\mathrm{kl}(\frac{1+\gamma}2\|\frac12)=\frac{\gamma^2}2(1+O(\gamma^2))$.
\end{proof}

\begin{remark}[What the adversary exploits, and how training can avoid it]\label{rem:decoding-training}
(a)~The worst case puts the whole error budget on small-margin states. Replace the $L^1$ budget by
$\E_\mu\max_t(e^\star_t)^r\le\eps^r$, and assume $\mu[\min_t\gamma_t\le u]\le Cu^\alpha$. Then \cref{thm:greedy}, applied to
$\mathsf D^r$, gives the exponent $\frac{r\alpha}{r+\alpha}$ for TV errors. This interpolates between $\frac\alpha{1+\alpha}$ ($r=1$)
and $\alpha$ ($r=\infty$). If the errors are independent of the margins, then
$\mu[\text{flip at }t]\le\E F_t(2e^\star_t)$, where $F_t(u)=\mu[\gamma_t\le u]$. For $\alpha\le1$, Jensen's inequality bounds
this by $C(2\E e^\star_t)^\alpha$, which is linear in the error when $\alpha=1$.
(b)~The factor $\pi^\star_t$ in \cref{thm:greedy-seq} is the price of training on sampled prefixes. Suppose that
$\mu=\cD$ and that, given each $x$, the training prefix law puts mass at least $\lambda$ on $s^\star_t(x)$. Examples
are $\nu_t=(1-\lambda)\dteach t+\lambda\,\delta_{s^\star_t}$ and soft-label distillation on the teacher's greedy outputs.
Then the training loss is at least $\lambda\eps^\star$, so \cref{thm:greedy} applies with $\eps^\star\le F/\lambda$, and no
path weight appears.
\end{remark}

\subsection{Hard labels on decoded outputs}

\begin{proposition}[A margin-free bound for sequence-level distillation]\label{prop:hard}
\Pv\Ck Fix a beam width $k\ge1$, with ties broken as in \cref{sec:beam}, so that $k=1$ is greedy decoding. Let
$\hat y(x)$ and $\hat y_\theta(x)$ be the teacher's and the student's outputs. If $\Qseq(\hat y(x))>\frac12$, then $\hat y_\theta(x)=\hat y(x)$. Hence
\[
\P_\mu\big[\hat y_\theta\ne\hat y\big]\le\E_\mu\min\Big\{1,\log_2\frac1{\Qseq(\hat y(x))}\Big\}\le\frac1{\ln2}\,\E_\mu\big[-\log\Qseq(\hat y(x))\big].
\]
For greedy decoding, a disagreement needs $q_t(a^\star_t\mid s^\star_t)\le\frac12$ at the first flip. The constant $\frac12$
cannot be improved.
\end{proposition}
\begin{proof}
If $Q(\hat y)>\frac12$, every prefix $\hat y_{\le t}$ has $Q$-mass above $\frac12$, so it has more mass than any other
prefix of the same length. By induction on $t$, $\hat y_{\le t}$ is a candidate at level $t$ and is the
student's best candidate, so it stays in the student's beam. At the end it is the student's most
probable hypothesis. Next, $\ind\{Q\le\frac12\}\le\min\{1,\log_2(1/Q)\}$. At a greedy flip,
$q(a^\star_t)\le q(b)\le1-q(a^\star_t)$. Finally, $q_t(\cdot\mid s^\star_t)=(\frac12-\delta,\frac12+\delta)$ on two tokens is a flip
for every $\delta>0$.
\end{proof}

\noindent
The right-hand side is the objective of sequence-level knowledge distillation \citep{kim2016sequence}:
maximum likelihood on the teacher's beam-search outputs. \Cref{prop:hard} shows that this objective is
a calibrated surrogate for decoding-level losslessness, with no margin in the constant. It asks the student
to be \emph{confident} on the teacher's decoded outputs, not to be faithful to the teacher's
distribution. These are different targets (\cref{cor:separation}). A soft-label student that matches a
teacher with small margins has $Q(\hat y)\approx P(\hat y)$, which may be far below $\frac12$. For such teachers
the margin-dependent bounds of \cref{thm:greedy,thm:greedy-seq} are the relevant ones.

\subsection{Beam search}\label{sec:beam}

Beam search of width $k$ keeps a set $\mathcal B_t$ of at most $k$ prefixes of length $t$, starting from
$\mathcal B_0=\{x\}$. The candidates at level $t$ are $\mathcal C_t=\{(u,a):u\in\mathcal B_{t-1},\ a\in\cV\}$. $\mathcal B_t$ consists of
the $k$ candidates of largest prefix probability under the decoding model. Ties are broken by a fixed
order on prefixes that agrees with the order on $\cV$ for candidates with a common parent, so that $k=1$
is greedy decoding. The output is the most probable element of $\mathcal B_T$. With the absorbing end-of-sequence
convention, every hypothesis has length $T$, so this is beam search without length normalisation. The
\emph{trajectory} is $(\mathcal B_1,\dots,\mathcal B_T,\text{output})$. For the teacher's beam search, let $P^+_t\ge P^-_t$ be
the $k$-th and $(k+1)$-th largest values of $\Pseq$ on $\mathcal C_t$, where $\Pseq(u)$ is the teacher's probability of
the prefix $u$. Let $P^+_{T+1}\ge P^-_{T+1}$ be the two largest values of $\Pseq$ on $\mathcal B_T$. A level with no
$(k+1)$-th candidate cannot change and is dropped, as is level $T+1$ when $k=1$. We write $P^x_{\le t}$ for
the law of $y_{\le t}$, and set $P^x_{\le T+1}:=\Pseq$.

\begin{theorem}[Beam search]\label{thm:beam}
\Pv\Ck Fix $x$ and $k\ge1$.
\begin{enumerate}[label=(\roman*),nosep]
\item If the student's trajectory first differs from the teacher's at level $t\in\{1,\dots,T+1\}$, then
$D_f(P^x_{\le t}\|Q^x_{\theta,\le t})\ge k_f(P^+_t,P^-_t)$ for every $f$-divergence. In particular
$\TV(P^x_{\le t},Q^x_{\theta,\le t})\ge\frac12(P^+_t-P^-_t)$, and
$\sum_{j\le\min(t,T)}\E_{s\sim\dteach j\mid x}\KL(p_j\|q_j)(s)=\KL(P^x_{\le t}\|Q^x_{\theta,\le t})\ge k_{\KL}(P^+_t,P^-_t)$.
\item (Beam-local form.) For $\mathsf D\in\{\TV,\KL\}$, the same lower bounds hold for
\[
\sum_{j\le\min(t,T)}\ \sum_{u\in\mathcal B_{j-1}}\Pseq(u)\,\mathsf D(p_j\|q_j)(u).
\]
Only per-token errors at the teacher's beam states enter, weighted by the teacher's probabilities of
those states.
\item (Exactness.) Let $\omega^{(k)}_f(x):=\min_tk_f(P^+_t,P^-_t)$. The outputs can differ only if
$D_f(\Pseq\|\Qseq)\ge\omega^{(k)}_f(x)$, so $\P_\mu[\text{beam outputs differ}]\le\bar\Phi_{\omega^{(k)}_f}(\cL_{D_f})$. If $f(0)<\infty$, or
if every candidate has positive teacher probability, then $\omega^{(k)}_f(x)$ is the infimum of $D_f(\Pseq\|\Qseq)$ over
students whose trajectory differs from the teacher's, and for trajectories the bound is sharp in the sense
of \cref{thm:greedy}(iii). All the $f$-divergences of \cref{prop:flip} except reverse KL have $f(0)<\infty$. For
reverse KL the condition is needed: with $T=1$, $p=(0.9,0.1,0,0)$ and $k=3$, the cut lies between two null
candidates, so $\omega^{(3)}_{\mathrm{rKL}}=0$, yet a finite $\KL(\Qseq\|\Pseq)$ forces $Q(c)=Q(d)=0$, and changing the beam then
costs at least $\log\frac1{0.9}$. For $k=1$, $(P^+_t,P^-_t)=\pi^\star_t\,(p^\star_{t,1},p^\star_{t,2})$, and (iii) is
\cref{thm:greedy-seq}, which needs no condition.
\item (The cut gap cannot be replaced by the output's own gaps.) Let $k=2$, $T=2$ and $\cV=\{a,b,c\}$. Take
first-token probabilities $(0.36,0.33,0.31)$ and prefix probabilities $aa:0.30$, $ab:0.059$, $ac:0.001$,
$ba:0.20$, $bb:0.10$, $bc:0.03$, $ca:0.309$ and $cb=cc:0.0005$. The teacher keeps $\{a,b\}$, then $\{aa,ba\}$, and outputs
$aa$. The cut gaps are $0.02$, $0.10$ and $0.10$, so $\omega^{(2)}_{\TV}=0.01$. The gaps between the ancestors of $aa$
and their best competitors are $P(a)-P(c)=0.05$, $P(aa)-P(bb)=0.20$ and $P(aa)-P(ba)=0.10$, which would give
the threshold $0.025$. Yet moving mass $0.01+\delta$ from the subtree of $b$ to that of $c$ makes the student keep
$\{a,c\}$, then $\{ca,aa\}$, and output $ca$, at $\TV=0.01+\delta$. The student ``repairs'' a search error of the
teacher, since $P(ca)=0.309>P(aa)$.
\end{enumerate}
\end{theorem}
\begin{proof}
(i) Let $t\le T$. The beams before level $t$ coincide, so the candidate sets $\mathcal C_t$ coincide. If every member
of the teacher's beam had larger $Q$-mass than every non-member, the student would keep the same beam,
whatever its tie rule. So $Q(a)\le Q(b)$ for some member $a$ and non-member $b$. Apply \cref{lem:merge} to the
law of $y_{\le t}$ with $S_+=\mathcal B_t$ and $S_-=\mathcal C_t\setminus\mathcal B_t$. For $t=T+1$, apply it with $S_+$ the teacher's
output and $S_-$ the rest of $\mathcal B_T$. The KL identity is \cref{thm:chain} applied to the first $t$ tokens.
(ii) Map $y_{\le t}$ to its path through the teacher's beam tree: to $y_{\le t}$ itself if $y_{<t}\in\mathcal B_{t-1}$, and
otherwise to $(j,y_{<j})$, where $j$ is the first level with $y_{\le j}\notin\mathcal B_j$. Candidates remain atoms, so
(i) holds for the image laws. These are the laws of a sequential process on the beam tree. Its step from
$u\in\mathcal B_{j-1}$ is $p_j(\cdot\mid u)$, with the non-beam children merged when $j<t$. The chain rule and the hybrid
argument of \cref{thm:chain,prop:hybrid} hold for such processes, with visiting probabilities $\Pseq(u)$, and
merging children can only decrease KL and TV.
(iii) The lower bound follows from (i), because $D_f(P\|Q)\ge D_f(P_{\le t}\|Q_{\le t})$ and the outputs can differ
only if the trajectories do. For attainment, take the law of \cref{lem:merge} at a minimising level and
extend it as in the proof of \cref{thm:greedy-seq}. The condition makes the strict clause of
\cref{lem:merge} available. If the trajectory has not changed before that level, the candidate sets
agree there, and the strict version reverses the $k$-th and $(k+1)$-th candidates.
(iv) Direct computation (\texttt{check\_decoding.py}).
\end{proof}

\subsection{Speculative decoding}

\begin{proposition}[Greedy versus sampled verification]\label{prop:spec}
\Pv\Ck Let the student draft tokens that the teacher verifies.
\begin{enumerate}[label=(\roman*),nosep]
\item (Speculative sampling.) A token drafted at state $s$ is accepted with probability
$\sum_a\min\{p_t(a\mid s),q_t(a\mid s)\}=1-\TV_t(s)$, and the output has exactly the teacher's law
\citep{leviathan2023fast}.
\item (Greedy verification.) Suppose the student drafts greedily and the teacher accepts the longest drafted
prefix that agrees with its own greedy continuation. This is the verification rule of blockwise parallel
decoding \citep{stern2018blockwise}, and the zero-temperature case of (i). Every verified prefix lies on the
teacher's greedy path. A round started at $s^\star_t$ with $L$ drafted tokens accepts
$N=\min\{j\ge0:\text{the student flips at }s^\star_{t+j}\}\wedge L$ of them. So
$\E_\mu N=\sum_{j=1}^L\big(1-\P_\mu[\text{a flip at one of }s^\star_t,\dots,s^\star_{t+j-1}]\big)$, and
\cref{thm:greedy,thm:greedy-seq}, applied to these windows, bound the expected acceptance length from below.
\item (Pointwise comparison.) At a state $s$ with margin $\gamma(s)>0$, the greedy rejection indicator is at most
$\min\{1,2\TV_t(s)/\gamma(s)\}$. This is $2/\gamma(s)$ times the rejection probability of speculative
sampling, and the constant $2$ cannot be improved.
\end{enumerate}
\end{proposition}
\begin{proof}
(i) The acceptance probability is $\sum_aq(a)\min\{1,p(a)/q(a)\}$. For exactness of the output law, see
\citet{leviathan2023fast}. (ii) The teacher emits its greedy token at each position, whether it accepts the
draft or corrects it. So each round starts at a teacher greedy state, and by \cref{lem:firstflip} the
accepted run ends at the first flip. (iii) This is the TV row of \cref{prop:flip}(ii).
\end{proof}

\noindent
Sampled verification rejects in proportion to the per-token TV, whatever the margins. Greedy verification
is a threshold rule. It never rejects while $\TV_t<\gamma_t/2$ along the greedy path. But when errors concentrate on
small-margin states, its worst-case rejection rate is of order $\eps^{\alpha/(1+\alpha)}$ (\cref{cor:rates}). This is a square
root when the margins have a bounded density at $0$, and it exceeds the linear rate of speculative sampling
for small $\eps$.

\begin{remark}[Measuring the bounds]\label{rem:measure-decoding}
All the quantities are cheap to compute. One greedy (or beam) decode of the teacher per input gives the
states $s^\star_t$, the margins $\gamma_t$, the top-two masses $m_t$ and the weights $\pi^\star_t$, or the cut probabilities
$P^\pm_t$. One student forward pass on the same prefixes gives the $e^\star_t$. The observed $\Ldec$ can then be
compared with $\bar\Phi_{\bar\kappa}(\eps^\star)$, with its per-step version, and with the path-weighted bound of
\cref{thm:greedy-seq} at the measured $T\eoff^{\KL}$. The ratio of the observed $\Ldec$ to the worst case shows how
strongly the student's errors concentrate on small-margin states.
\end{remark}
